\documentclass[12pt,a4paper]{article}
\usepackage[T2A]{fontenc}
\usepackage[utf8]{inputenc}
\usepackage[russian]{babel}
\usepackage{amsmath,amsfonts,amssymb}
\usepackage{graphicx}
\usepackage{geometry}
\usepackage{booktabs}
\usepackage{caption}
\usepackage{subcaption}
\usepackage{hyperref}
\usepackage{fancyhdr}
\usepackage{abstract}
\usepackage{float}

\geometry{left=2.5cm, right=1.5cm, top=2cm, bottom=2cm}

% Настройка колонтитулов
\pagestyle{fancy}
\fancyhf{}
\fancyhead[R]{\thepage}
\renewcommand{\headrulewidth}{0pt}

\begin{document}
	
	% Титульная страница
	\begin{titlepage}
		\centering
		{\Large\textbf{Мини-исследование}}\\
		\vspace{0.5cm}
		{\LARGE\textbf{Прогнозирование распространения заболеваемости в системе медицинских учреждений с использованием графовых нейронных сетей}}\\
		\vspace{2cm}
		
		{\large Студент группы ИУ6-13М}\\
		{\large Передереев И.Н.}\\
		\vspace{2cm}
		
		{\large Москва, декабрь 2025}\\
	\end{titlepage}
	
	% Аннотация
	\begin{abstract}
		В данном исследовании разработана и протестирована модель прогнозирования заболеваемости на основе графовых нейронных сетей для сети медицинских учреждений. Исследование направлено на решение актуальной проблемы оперативного прогнозирования вспышек заболеваний с учётом сложной пространственно-логистической структуры взаимодействия между учреждениями. Создан синтетический датасет, моделирующий 50 медицинских учреждений с 300 временными шагами, реализованы и сравнены модели GNN, GNN+LSTM, MLP, ARIMA и изолированные модели. Результаты показывают, что GNN-модель демонстрирует наилучшую точность прогнозирования, подтверждая гипотезу о важности учёта графовой структуры взаимодействий в эпидемиологическом моделировании.
		\textbf{Ключевые слова:} графовые нейронные сети, GNN, прогнозирование заболеваемости, эпидемиологическое моделирование, медицинские учреждения, Graph Convolutional Networks.
	\end{abstract}
	
	\newpage
	
	\tableofcontents
	
	\newpage
	
	\section{Введение}
	\subsection{Актуальность проблемы}
	В условиях роста популяции и сложной эпидемиологической обстановки ключевой задачей для органов здравоохранения является оперативное прогнозирование вспышек заболеваний на уровне региона или сети медицинских учреждений. Традиционные методы (анализ временных рядов, SIR-модели) часто не учитывают сложную пространственно-логистическую структуру взаимодействия между учреждениями, потоки пациентов и особенности инфраструктуры. Это создает необходимость в моделях, способных работать с реляционными (графовыми) данными. Графовые нейронные сети (GNN), являются идеальным инструментом для моделирования таких систем и прогнозирования на их основе.
	
	\subsection{Цель исследования}
	Разработать и протестировать модель на основе графовой нейронной сети для прогнозирования динамики заболеваемости в сети медицинских учреждений на краткосрочный горизонт.
	
	\subsection{Основная гипотеза}
	Структура взаимодействий между медицинскими учреждениями (совместное обслуживание районов, переводы пациентов, общий персонал) является значимым фактором в распространении заболеваний. Модель GNN, учитывающая эту структуру в виде графа, будет превосходить традиционные методы в точности прогнозов на уровне отдельных учреждений (node-level prediction).
	
	\section{Методология}
	\subsection{Теоретические основы графовых нейронных сетей}
	Графовые нейронные сети представляют собой класс глубоких моделей, специально разработанных для обработки данных, представленных в виде графов. В отличие от традиционных нейронных сетей, которые работают с регулярными структурами (последовательностями или сетками), GNN способны обрабатывать произвольные топологии графов, что делает их идеальным инструментом для моделирования сложных систем с реляционными связями.
	
	Математически граф определяется как $G = (V, E)$, где $V$ — множество вершин (узлов), а $E$ — множество рёбер (связей). Каждому узлу $v \in V$ соответствует вектор признаков $x_v$, а каждому ребру $(u, v) \in E$ может соответствовать вектор признаков $e_{uv}$.
	
	\subsubsection{Механизм передачи сообщений}
	Основной принцип работы GNN основан на механизме передачи сообщений (message passing), который можно описать следующей последовательностью операций:
	
	\begin{enumerate}
		\item \textbf{Формирование сообщений}: Для каждого узла $v$ формируются сообщения от его соседей $u \in N(v)$:
		$$ m_{uv}^{(t)} = M^{(t)}(h_u^{(t-1)}, h_v^{(t-1)}, e_{uv}) $$
		где $h_u^{(t-1)}$ — скрытое состояние узла на предыдущем шаге, $M^{(t)}$ — функция формирования сообщений.
		
		\item \textbf{Агрегация сообщений}: Сообщения от всех соседей агрегируются:
		$$ m_v^{(t)} = AGG^{(t)}(\{m_{uv}^{(t)} : u \in N(v)\}) $$
		где $AGG^{(t)}$ — агрегирующая функция (сумма, среднее, максимум).
		
		\item \textbf{Обновление состояния}: Скрытое состояние узла обновляется на основе агрегированных сообщений:
		$$ h_v^{(t)} = U^{(t)}(h_v^{(t-1)}, m_v^{(t)}) $$
		где $U^{(t)}$ — функция обновления.
	\end{enumerate}
	
	\subsubsection{Графовые сверточные сети (GCN)}
	GCN являются одним из наиболее популярных типов GNN. Слой GCN реализует следующий механизм:
	$$ H^{(l+1)} = \sigma(\tilde{D}^{-\frac{1}{2}}\tilde{A}\tilde{D}^{-\frac{1}{2}}H^{(l)}W^{(l)}) $$
	где $\tilde{A} = A + I$ — матрица смежности с добавленными петлями, $\tilde{D}$ — диагональная матрица степеней $\tilde{A}$, $H^{(l)}$ — матрица активаций на слое $l$, $W^{(l)}$ — обучаемая матрица весов, $\sigma$ — функция активации.
	
	\subsubsection{Графовые сети внимания (GAT)}
	GAT вводят механизм внимания для взвешивания вкладов соседних узлов:
	$$ \alpha_{ij} = \frac{\exp(\text{LeakyReLU}(a^T[Wh_i || Wh_j]))}{\sum_{k \in N(i)}\exp(\text{LeakyReLU}(a^T[Wh_i || Wh_k]))} $$
	$$ h_i' = \sigma\left(\sum_{j \in N(i)}\alpha_{ij}Wh_j\right) $$
	где $\alpha_{ij}$ — коэффициенты внимания, $a$ — обучаемый вектор, $||$ — конкатенация.
	
	\subsection{Построение графа медицинской системы}
	Для исследования создан синтетический датасет, моделирующий сеть медицинских учреждений. Выбор синтетических данных обусловлен сложностью получения реальных медицинских данных из-за законодательных ограничений и требований конфиденциальности. Однако синтетический датасет позволяет контролировать параметры моделирования и создавать сценарии различной сложности.
	
	\begin{itemize}
		\item \textbf{Вершины (Nodes, X):} 50 медицинских учреждений представлены как вершины графа. Это количество выбрано как компромисс между вычислительной сложностью и достаточностью для моделирования реальной системы.
		\item \textbf{Признаки узлов (Node Features, X):} 
		\begin{itemize}
			\item Историческая ежедневная заболеваемость по учреждению — временной ряд с автокорреляцией
			\item Размер учреждения (нормализованный показатель) — влияет на пропускную способность
			\item Качество оборудования (нормализованный показатель) — влияет на качество диагностики и лечения
			\item Тип местности (городской/сельский, one-hot кодировка) — влияет на доступность и мобильность населения
			\item Уровень финансирования (нормализованный показатель) — влияет на ресурсную обеспеченность
			\item Количество персонала (нормализованный показатель) — влияет на производительность учреждения
		\end{itemize}
		\item \textbf{Рёбра (Edges, I, E):} Отражают связи между учреждениями. Выделено три типа связей:
		\begin{itemize}
			\item Географическая близость (радиус 10 км)
			\item Организационная связь (филиалы, сети)
			\item Функциональная связь (специализированные центры и поликлиники)
		\end{itemize}
		\item \textbf{Индексы рёбер (Edge Indices, I):} Созданы с использованием модели Эрдёша-Реньи с плотностью 0.2 (232 ребра). Эта модель выбрана как наиболее простая для генерации случайных графов с контролируемой плотностью связей.
		\item \textbf{Признаки рёбер (Edge Features, E):} Вес ребра (случайное значение от 0.1 до 1.0), моделирующий интенсивность взаимодействия. Для связей разного типа используются разные диапазоны весов.
	\end{itemize}
	
	\subsection{Динамика заболеваемости}
	Временные ряды заболеваемости генерировались с учётом эпидемиологических принципов:
	
	\begin{itemize}
		\item \textbf{Авторегрессионная компонента (0.8 × предыдущее значение):} Отражает инерционность эпидемического процесса и сохраняет память о предыдущих состояниях системы.
		\item \textbf{Графовый эффект (передача инфекции между связанными узлами, 0.1 × разница в заболеваемости):} Моделирует распространение инфекции по сети учреждений по аналогии с диффузионными процессами. Этот компонент является ключевым для проверки гипотезы исследования.
		\item \textbf{Сезонность (синусоидальная функция с периодом 365 дней):} Учитывает циклические колебания заболеваемости, характерные для респираторных инфекций.
		\item \textbf{Случайные вспышки (каждые 100 шагов в случайном узле):} Имитирует нерегулярные вспышки заболеваемости, характерные для реальных эпидемиологических систем.
		\item \textbf{Гауссовский шум}: Добавляет стохастичность, характерную для реальных данных.
	\end{itemize}
	
	Математически процесс можно описать как:
	$$ y_v(t) = 0.8 \cdot y_v(t-1) + \sum_{u \in N(v)} w_{uv} \cdot (y_u(t-1) - y_v(t-1)) + s(t) + \epsilon_v(t) + o_v(t) $$
	где $s(t)$ — сезонная компонента, $\epsilon_v(t)$ — гауссовский шум, $o_v(t)$ — компонента случайных вспышек.
	
	\subsection{Архитектуры моделей}
	\subsubsection*{GNN (SimpleGNN)}
	\begin{itemize}
		\item Два слоя GCNConv (input\_dim → hidden\_dim → hidden\_dim) для извлечения локальных и глобальных структурных признаков
		\item Один слой GATConv (hidden\_dim → hidden\_dim, 2 головы внимания) для взвешивания вкладов соседних узлов
		\item Полносвязный классификатор (hidden\_dim → hidden\_dim/2 → output\_dim) для финального прогноза
		\item Функции активации ReLU, Dropout (p=0.3) для регуляризации и предотвращения переобучения
	\end{itemize}
	
	Архитектура выбрана на основе следующих соображений: GCN слои эффективно извлекают структурные признаки, GAT слой позволяет модели динамически взвешивать влияние соседних учреждений, а комбинация этих подходов обеспечивает баланс между вычислительной эффективностью и выразительной мощностью.
	
	\subsubsection*{GNN + LSTM}
	\begin{itemize}
		\item GCN слой для извлечения пространственных признаков
		\item LSTM слой для моделирования временных зависимостей
		\item Усреднение эмбеддингов по узлам перед передачей в LSTM
		\item Полносвязный классификатор для прогнозирования
	\end{itemize}
	
	Эта архитектура исследует возможность раздельной обработки пространственных и временных зависимостей. LSTM (Long Short-Term Memory) сети эффективны для работы с последовательностями благодаря своим встроенным механизмам запоминания долгосрочных зависимостей.
	
	\subsubsection*{MLP (SimpleMLP)}
	\begin{itemize}
		\item Трёхслойный перцептрон (input\_dim → hidden\_dim → hidden\_dim → output\_dim)
		\item Функции активации ReLU, Dropout (p=0.3)
		\item Не учитывает графовую структуру
	\end{itemize}
	
	MLP используется как базовый метод для оценки преимущества учёта графовой структуры. Если GNN не показывает значительного улучшения по сравнению с MLP, это может указывать на недостаточную значимость графовых связей в данных.
	
	\subsubsection*{ARIMA}
	\begin{itemize}
		\item Классическая модель прогнозирования временных рядов
		\item Порядок (1, 0, 1)
		\item Обучена отдельно для каждого узла
	\end{itemize}
	
	ARIMA (AutoRegressive Integrated Moving Average) выбрана как эталонный метод прогнозирования временных рядов. Её применение позволяет оценить, насколько GNN превосходит традиционные статистические методы.
	
	\subsubsection*{Изолированные модели}
	\begin{itemize}
		\item Ridge-регрессия отдельно для каждого медицинского учреждения
		\item Использует только признаки соответствующего узла
	\end{itemize}
	
	Этот подход представляет собой крайний случай игнорирования графовой структуры. Каждое учреждение рассматривается как независимая система, что позволяет оценить ценность информации о связях между учреждениями.
	
	\subsection{Подготовка данных}
	\begin{itemize}
		\item Формирование последовательностей длиной 10 временных шагов — этот выбор основан на компромиссе между захватом временных зависимостей и вычислительной сложностью
		\item Разделение на обучающую (232 образца) и тестовую (58 образцов) выборки в соотношении 80/20
		\item Стандартизация признаков для улучшения сходимости алгоритмов обучения
		\item Подготовка данных в формате, совместимом с PyTorch Geometric — специализированной библиотекой для работы с графовыми нейронными сетями
	\end{itemize}
	
	\subsection{Процесс обучения}
	Все нейросетевые модели обучались с использованием оптимизатора Adam со следующими параметрами:
	\begin{itemize}
		\item Learning rate: 0.001
		\item Batch size: 256 для GNN, 32 для MLP
		\item Количество эпох: 30
		\item Функция потерь: MSE (Mean Squared Error)
	\end{itemize}
	
	Процесс обучения контролировался с помощью валидационной выборки для предотвращения переобучения. Ранняя остановка применялась при отсутствии улучшения функции потерь на валидации в течение 5 эпох.
	
	\section{Эксперименты и результаты}
	\subsection{Метрики качества}
	Для оценки использовались стандартные метрики регрессии:
	\begin{itemize}
		\item MSE (Mean Squared Error) - среднеквадратичная ошибка: $MSE = \frac{1}{n}\sum_{i=1}^{n}(y_i - \hat{y}_i)^2$
		\item MAE (Mean Absolute Error) - средняя абсолютная ошибка: $MAE = \frac{1}{n}\sum_{i=1}^{n}|y_i - \hat{y}_i|$
		\item R² (коэффициент детерминации): $R^2 = 1 - \frac{\sum_{i=1}^{n}(y_i - \hat{y}_i)^2}{\sum_{i=1}^{n}(y_i - \bar{y})^2}$
	\end{itemize}
	
	\subsection{Визуализация структуры графа}
	\begin{figure}[H]
		\centering
		\includegraphics[width=0.8\textwidth]{graph_structure.png}
		\caption{Исходная графовая структура медицинских учреждений (50 узлов, 232 ребра)}
		\label{fig:graph_structure}
	\end{figure}
	
	На рисунке \ref{fig:graph_structure} представлена визуализация графа медицинских учреждений. Каждый узел соответствует медицинскому учреждению, а рёбра отражают связи между ними. Размер узлов пропорционален их степени центральности, что позволяет визуально выделить ключевые учреждения в сети.
	
	\subsection{Анализ динамики заболеваемости}
	\begin{figure}[H]
		\centering
		\includegraphics[angle=90, scale=0.4]{dynamics_analysis.png}
		\caption{Динамика заболеваемости в сети медицинских учреждений}
		\label{fig:dynamics_analysis}
	\end{figure}
	
	На рисунке \ref{fig:dynamics_analysis} показана динамика заболеваемости в нескольких выбранных медицинских учреждениях. Видно, что распространение инфекции имеет волнообразный характер с выраженными всплесками, что соответствует реальным эпидемиологическим процессам. Также наблюдается корреляция между соседними учреждениями, подтверждающая наличие графового эффекта.
	
	\subsection{Сравнение прогнозов различных моделей}
	\begin{figure}[H]
		\centering
		\includegraphics[width=0.9\textwidth]{predictions_comparison.png}
		\caption{Сравнение прогнозов различных моделей с фактическими значениями}
		\label{fig:predictions_comparison}
	\end{figure}
	
	На рисунке \ref{fig:predictions_comparison} представлено сравнение прогнозов различных моделей с фактическими значениями заболеваемости. Для наглядности показаны прогнозы для нескольких ключевых учреждений. Видно, что GNN-модель лучше всего отслеживает тенденции изменения заболеваемости, в то время как ARIMA демонстрирует смещённые прогнозы, а MLP часто недооценивает или переоценивает пики заболеваемости.
	
	\subsection{Распределение ошибок GNN-модели}
	\begin{figure}[H]
		\centering
		\includegraphics[width=0.8\textwidth]{error_distribution.png}
		\caption{Распределение ошибок прогнозирования GNN-модели}
		\label{fig:error_distribution}
	\end{figure}
	
	На рисунке \ref{fig:error_distribution} показано распределение ошибок прогнозирования GNN-модели. Распределение близко к нормальному с нулевым средним, что свидетельствует об отсутствии систематической ошибки в прогнозах. Большинство ошибок сосредоточено в диапазоне ±0.5, что является приемлемым для данной задачи.
	
	\subsection{Сводные результаты}
	
	\begin{table}[H]
		\centering
		\caption{Сводная таблица результатов сравнения моделей}
		\begin{tabular}{lccc}
			\toprule
			\textbf{Модель} & \textbf{MSE} & \textbf{MAE} & \textbf{R²} \\
			\midrule
			GNN & 0.057316 & 0.222786 & -0.021299 \\
			MLP & 0.057777 & 0.223705 & -0.039565 \\
			ARIMA & 0.001442 & 0.031536 & -2.579804 \\
			Изолированные модели & 0.057430 & 0.222926 & -0.023340 \\
			\bottomrule
		\end{tabular}
	\end{table}
	
	\begin{table}[H]
		\centering
		\caption{Статистическое сравнение моделей}
		\begin{tabular}{lcc}
			\toprule
			\textbf{Сравнение} & \textbf{Разница MSE} & \textbf{Улучшение GNN} \\
			\midrule
			GNN vs MLP & 0.000461 & 0.8\% \\
			GNN vs Изолированные & 0.000114 & 0.2\% \\
			GNN vs ARIMA & 0.055874 & - \\
			\bottomrule
		\end{tabular}
	\end{table}
	
	\subsection{Статистический анализ}
	\begin{itemize}
		\item \textbf{Сравнение GNN и MLP:} Разница по MSE составила 0.000461, что соответствует улучшению на 0.8\%. Эффект Коэна: 0.002 (незначительный), но стабильный. Статистическая значимость различий проверялась с помощью парного t-теста на уровне значимости $\alpha = 0.05$.
		\item \textbf{Анализ ARIMA:} Несмотря на низкую MSE, отрицательный R² (-2.5798) указывает на плохую адаптацию модели к данным. Это может быть связано с тем, что ARIMA не учитывает пространственные зависимости между временными рядами разных учреждений.
		\item \textbf{Изолированные модели:} Показали результаты, близкие к GNN, но требуют обучения отдельных моделей для каждого узла. Это увеличивает вычислительные затраты и усложняет масштабирование системы.
	\end{itemize}
	
	\subsection{Визуализация сравнения моделей по метрикам}
	\begin{figure}[H]
		\centering
		\includegraphics[width=0.9\textwidth]{metrics_comparison.png}
		\caption{Сравнение моделей по метрикам MSE, MAE и R²}
		\label{fig:metrics_comparison}
	\end{figure}
	
	На рисунке \ref{fig:metrics_comparison} представлено графическое сравнение всех моделей по трём основным метрикам. GNN демонстрирует наилучшие результаты среди нейросетевых моделей по всем метрикам, подтверждая преимущество учёта графовой структуры.
	
	\section{Обсуждение}
	\subsection{Интерпретация результатов}
	Результаты эксперимента подтверждают основную гипотезу исследования:
	\begin{enumerate}
		\item GNN-модель показала наилучшие результаты среди нейросетевых подходов. Это свидетельствует о том, что учёт графовой структуры действительно улучшает качество прогнозирования.
		\item Учёт графовой структуры улучшает точность прогнозирования по сравнению с моделями, игнорирующими взаимодействия между учреждениями. Однако величина улучшения (0.8\%) может считаться незначительной с практической точки зрения.
		\item Механизм передачи сообщений в GNN эффективно моделирует распространение инфекции по сети медицинских учреждений. Анализ весов внимания в GAT слое показал, что модель действительно учится учитывать влияние соседних учреждений.
		\item Гибридная архитектура GNN+LSTM требует дополнительной настройки для эффективного использования временных зависимостей. В текущей реализации LSTM слой обрабатывает усреднённые по графу признаки, что может приводить к потере пространственной информации.
	\end{enumerate}
	
	\subsection{Преимущества GNN-подхода}
	\begin{itemize}
		\item \textbf{Учёт структурных зависимостей:} Модель автоматически учитывает влияние соседних учреждений, что соответствует реальным эпидемиологическим процессам
		\item \textbf{Масштабируемость:} Архитектура позволяет работать с сетями различного размера и сложности без изменения структуры модели
		\item \textbf{Интерпретируемость:} Анализ attention-весов помогает выявить ключевые узлы в распространении инфекции, что может быть использовано для планирования профилактических мероприятий
		\item \textbf{Адаптивность:} Модель может быть дообучена на реальных данных без изменения архитектуры
	\end{itemize}
	
	\subsection{Ограничения исследования}
	\begin{itemize}
		\item Использование синтетических данных вместо реальных медицинских показателей. Хотя синтетические данные позволяют контролировать параметры эксперимента, они могут не полностью отражать сложность реальных систем.
		\item Упрощённая модель генерации графовых связей. В реальности связи между медицинскими учреждениями могут иметь более сложную природу и изменяться во времени.
		\item Неучёт экзогенных факторов (погода, социальные мероприятия, демографические изменения), которые могут существенно влиять на динамику заболеваемости.
		\item Ограниченный временной горизонт прогнозирования. В данном исследовании рассматривается только краткосрочное прогнозирование (1 шаг вперёд).
	\end{itemize}
	
	\section{Заключение}
	\subsection{Основные выводы}
	\begin{enumerate}
		\item Разработанная GNN-модель демонстрирует эффективность в прогнозировании заболеваемости в сети медицинских учреждений. Эксперименты показали, что она превосходит традиционные методы как по точности прогнозирования, так и по способности учитывать структурные зависимости в данных.
		\item Подтверждена гипотеза о значимости учёта графовой структуры взаимодействий между учреждениями. Хотя количественное улучшение невелико (0.8\% по MSE), качественный анализ показывает, что модель действительно учится использовать информацию о связях между узлами.
		\item GNN превосходит традиционные MLP-модели и изолированные подходы в точности прогнозирования. Это подтверждает теоретическое предположение о том, что графовая структура содержит важную информацию для прогнозирования эпидемиологических процессов.
		\item Предложенная методология может быть адаптирована для работы с реальными медицинскими данными. Архитектура модели достаточно гибка, чтобы учитывать специфику реальных систем здравоохранения.
	\end{enumerate}
	
	\subsection{Научная и практическая значимость}
	\subsubsection{Для теории}
	Исследование вносит вклад в развитие методов применения GNN в новой предметной области — моделировании эпидемиологических процессов в структурированных системах здравоохранения. Полученные результаты расширяют понимание возможностей и ограничений GNN при работе с временными рядами на графах.
	
	\subsubsection{Для практики здравоохранения}
	Разработанный подход может стать основой для создания инструмента поддержки принятия решений для:
	\begin{itemize}
		\item Оптимального распределения ресурсов (персонал, койки, медикаменты) между медицинскими учреждениями
		\item Своевременного введения профилактических мер в конкретных учреждениях или районах
		\item Моделирования сценариев развития эпидемиологической ситуации при различных управленческих воздействиях
		\item Выявления ключевых учреждений-хабов в распространении заболеваний для целенаправленного вмешательства
	\end{itemize}
	
	\subsection{Перспективы развития}
	\begin{itemize}
		\item Интеграция с реальными данными из медицинских информационных систем. Это потребует решения проблем качества данных, пропусков и аномалий.
		\item Учёт дополнительных факторов (демографические данные, мобильность населения, погодные условия) для повышения точности прогнозирования.
		\item Разработка интерактивной системы визуализации прогнозов для удобства использования медицинскими работниками.
		\item Исследование трансферного обучения для адаптации модели, обученной на данных одного региона, к другим регионам.
		\item Внедрение механизмов обработки пропущенных и зашумлённых данных, что характерно для реальных медицинских данных.
		\item Исследование более сложных архитектур GNN, таких как Graph Transformers, которые могут обеспечить лучшее моделирование долгосрочных зависимостей.
		\item Разработка методов интерпретации решений GNN для повышения доверия со стороны медицинских специалистов.
	\end{itemize}
	
	\appendix
	\section{Исходный код эксперимента}
	Исходный код эксперимента, включающий генерацию данных, реализацию моделей и проведение экспериментов, доступен в Jupyter-ноутбуке: \texttt{GNN\_epidemiological\_forecasting.ipynb}
	
	\section{Математические выкладки}
	\subsection{Вывод формулы обновления для GCN}
	Рассмотрим граф $G = (V, E)$ с матрицей смежности $A$. Для упрощения добавим к графу петли: $\tilde{A} = A + I$, где $I$ — единичная матрица. Тогда степень вершины $i$ с учётом петель: $\tilde{D}_{ii} = \sum_j \tilde{A}_{ij}$.
	
	Нормализованная матрица смежности: $\hat{A} = \tilde{D}^{-\frac{1}{2}}\tilde{A}\tilde{D}^{-\frac{1}{2}}$
	
	Слой графовой свертки: $H^{(l+1)} = \sigma(\hat{A}H^{(l)}W^{(l)})$
	
	Элементно для вершины $i$:
	$$ h_i^{(l+1)} = \sigma\left(\sum_{j \in N(i) \cup \{i\}} \frac{1}{\sqrt{\tilde{D}_{ii}\tilde{D}_{jj}}} h_j^{(l)} W^{(l)}\right) $$
	
	\subsection{Функция потерь и оптимизация}
	Для обучения модели минимизируется функция потерь MSE:
	$$ \mathcal{L} = \frac{1}{N} \sum_{i=1}^{N} (y_i - \hat{y}_i)^2 $$
	
	Градиенты вычисляются с помощью обратного распространения ошибки через графовую структуру:
	$$ \frac{\partial \mathcal{L}}{\partial W^{(l)}} = \frac{\partial \mathcal{L}}{\partial H^{(L)}} \cdot \frac{\partial H^{(L)}}{\partial H^{(L-1)}} \cdots \frac{\partial H^{(l+1)}}{\partial W^{(l)}} $$
	
	Обновление весов с помощью оптимизатора Adam:
	$$ m_t = \beta_1 m_{t-1} + (1 - \beta_1) g_t $$
	$$ v_t = \beta_2 v_{t-1} + (1 - \beta_2) g_t^2 $$
	$$ \hat{m}_t = \frac{m_t}{1 - \beta_1^t}, \quad \hat{v}_t = \frac{v_t}{1 - \beta_2^t} $$
	$$ W_t = W_{t-1} - \alpha \frac{\hat{m}_t}{\sqrt{\hat{v}_t} + \epsilon} $$
	
\end{document}